\section{Introdução}

O projeto Pauliceia 2.0 articula pesquisa histórica, informação geográfica e colaboração para o mapeamento da história de São Paulo. Ferla et al.~\cite{ferla2020} apresentam sua concepção como uma iniciativa de código aberto associada às humanidades digitais, à história pública e à ciência aberta. Nesse contexto, localizar um endereço envolve relacionar a expressão fornecida pelo pesquisador aos registros disponíveis e, posteriormente, consultar a informação espacial e temporal correspondente. A recuperação do nome é, portanto, uma etapa importante para o acesso ao acervo georreferenciado.

Os nomes presentes em um cadastro podem diferir da consulta por acentuação, abreviações, omissões, transposições e outras alterações de escrita. Há também grafias distintas registradas em fontes documentais. A igualdade literal é insuficiente para acomodar essas diferenças. Por outro lado, aceitar indiscriminadamente nomes parecidos pode aproximar logradouros distintos, sobretudo quando compartilham parte das palavras. O problema investigado consiste em ordenar os nomes candidatos para favorecer a referência esperada, sem assumir que semelhança textual comprova identidade geográfica ou histórica.

Na recuperação lexical, BM25 atribui escores aos documentos a partir da correspondência entre termos, com tratamento de frequência e comprimento~\cite{elasticSimilarity}. Essa representação pode perder evidência quando uma pequena alteração modifica um termo inteiro. A indexação adicional por trigramas permite compartilhar fragmentos de caracteres entre expressões diferentes. Contudo, recuperar a referência entre os candidatos e colocá-la na primeira posição são objetivos distintos: uma lista pode ter cobertura adequada e ainda exigir que o usuário procure a opção correta em posições pouco visíveis.

A combinação de similaridades já possui antecedentes na literatura. Santos, Murrieta-Flores e Martins~\cite{santos2018} avaliaram medidas textuais para correspondência de topônimos e sua combinação por aprendizado supervisionado, com resultados favoráveis a conjuntos de árvores de decisão. Esse antecedente fundamenta a abordagem, mas não determina seu desempenho no catálogo e nas consultas do Pauliceia. Nosso estudo investiga essa aplicação por meio de recuperação em duas etapas, avaliação por posição e integração a um portal existente. Não propõe a combinação supervisionada de medidas como uma técnica inédita.

A pergunta de pesquisa é: \emph{em que medida a reordenação supervisionada de candidatos recuperados por BM25 e trigramas melhora a recuperação de nomes de logradouros no recorte disponível do Pauliceia?} Para respondê-la, comparamos BM25, BM25 com trigramas e Random Forest em consultas separadas por grupos de nomes. Também comparamos reordenadores baseados em uma única similaridade e examinamos um conjunto complementar de grafias documentadas e transformações simuladas. A avaliação da integração verificou a correspondência das sugestões com o processamento offline, a entrega dos pontos cadastrados e a latência HTTP local.

As contribuições são a implementação reproduzível dessa arquitetura no protótipo local, a comparação empírica dos métodos no catálogo estudado e a explicitação dos limites entre recuperação textual, atendimento pelo portal e validação histórica. No teste reservado, a reordenação corrigiu as 36 falhas de primeira posição dos trigramas. Já no complemento, uma medida isolada empatou com o modelo supervisionado. Esses resultados motivam uma análise das condições em que a combinação ajuda, sem generalizar o desempenho observado para toda a documentação histórica ou para consultas reais ainda não avaliadas.
